\documentclass[11pt,a4paper]{article}

\usepackage[italian]{babel}
\usepackage[T1]{fontenc}
\usepackage[utf8]{inputenc}
\usepackage{mathpazo}
\usepackage{amsmath,amssymb}
\usepackage{geometry}
\usepackage{enumitem}
\usepackage{fancyhdr}
\usepackage{lastpage}
\usepackage{microtype}
\usepackage{xcolor}
\usepackage{tabularx}
\usepackage{booktabs}

\geometry{top=2.05cm,bottom=2.0cm,left=2.15cm,right=2.15cm}
\definecolor{examNavy}{HTML}{18324A}
\definecolor{examGold}{HTML}{AD8549}
\definecolor{examPale}{HTML}{F5F2EC}
\definecolor{examRule}{HTML}{D8D2C7}
\setlength{\parindent}{0pt}
\setlength{\parskip}{0.3em}
\setlength{\headheight}{15pt}
\pagestyle{fancy}
\fancyhf{}
\fancyhead[L]{\small\textsc{Corso di laurea in Informatica}}
\fancyhead[R]{\small\textsc{Matematica discreta e logica}}
\fancyfoot[C]{\small Pagina \thepage\ di \pageref{LastPage}}
\renewcommand{\headrulewidth}{0.5pt}
\renewcommand{\headrule}{\hbox to\headwidth{\color{examGold}\leaders\hrule height \headrulewidth\hfill}}
\renewcommand{\footrulewidth}{0pt}
\setlist[enumerate,1]{label=\textbf{(\Alph*)},leftmargin=1.35cm,itemsep=0.12em,topsep=0.18em}
\renewcommand{\arraystretch}{1.45}

\newcommand{\quesito}[2]{%
    \par\addvspace{0.55em}%
    \noindent\colorbox{examPale}{%
        \parbox{\dimexpr\linewidth-2\fboxsep\relax}{%
            \color{examNavy}\textbf{QUESITO #1}\hfill\textbf{#2 punti}}}%
    \par\nobreak\vspace{0.15em}%
}

\newcommand{\casella}{\fbox{\rule{0pt}{1.1ex}\hspace{1.1ex}}}

\begin{document}

\begin{center}
    \begin{tabularx}{\linewidth}{@{}>{\centering\arraybackslash}p{1.65cm}X@{}}
        \fcolorbox{examGold}{examPale}{%
            \parbox[c][1.55cm][c]{1.4cm}{\centering
                {\small\scshape Corso}\\[0.1em]
                {\Large\bfseries\color{examNavy} Informatica}}}
        &
        {\raggedright
            {\small\scshape Universit\`a degli Studi}\\[0.2em]
            {\LARGE\bfseries\color{examNavy} Matematica discreta e logica}\\[0.15em]
            {\large Prova scritta}}
    \end{tabularx}
    \vspace{0.45em}
    {\color{examGold}\rule{\linewidth}{1.1pt}}
\end{center}

\vspace{0.2em}
\begin{tabularx}{\linewidth}{@{}lXr@{}}
    \textbf{Durata} & 120 minuti & \textbf{Data} \rule{2.8cm}{0.4pt}\\[0.8em]
    \textbf{Cognome e nome} & \multicolumn{2}{l}{\rule{12cm}{0.4pt}}\\[0.8em]
    \textbf{Matricola} & \multicolumn{2}{l}{\rule{5cm}{0.4pt}}
\end{tabularx}

\vspace{0.6em}
\colorbox{examPale}{%
    \parbox{\dimexpr\linewidth-2\fboxsep\relax}{%
        \textbf{\color{examNavy}Indicazioni e punteggio}\\[-0.1em]
        Per ogni quesito selezionare una sola risposta.
        Dal quesito 1 al quesito 8 sono assegnati 3 punti per ogni risposta
        corretta. Dal quesito 9 al quesito 11 sono assegnati 2 punti per ogni
        risposta corretta. Alle risposte errate o omesse sono assegnati 0 punti.
        Il voto in trentesimi coincide con il punteggio totale.
        Punteggio massimo: \textbf{30/30}. Voto minimo: \textbf{18/30}.}}

\begin{samepage}
\quesito{1}{3}
La successione definita dalla relazione ricorsiva
\[
    a_{n+2}=12a_{n+1}-35a_n+2^n,
    \qquad a_0=1,\quad a_1=2,
\]
\`e
\begin{enumerate}
    \item $\displaystyle a_n=\frac{7\cdot5^n-7\cdot7^n+2^n}{15}$
    \item $\displaystyle a_n=\frac{35\cdot5^n+21\cdot7^n+2^n}{15}$
    \item $\displaystyle a_n=\frac{35\cdot5^n-21\cdot7^n+2^n}{15}$
    \item $\displaystyle a_n=\frac{35\cdot5^n-21\cdot7^n-2^n}{15}$
    \item $\displaystyle a_n=5^n+7^n+2^n$
\end{enumerate}
\end{samepage}

\begin{samepage}
\quesito{2}{3}
Il numero di funzioni suriettive $f:A\to B$, dove
$A=\{1,2,3,4,5,6\}$ e $B=\{a,b,c\}$, \`e
\begin{enumerate}
    \item $360$
    \item $540$
    \item $729$
    \item $720$
    \item $486$
\end{enumerate}
\end{samepage}

\begin{samepage}
\quesito{3}{3}
Il numero di interi $x$ tali che $500\leq x\leq1500$ e
\[
    \begin{cases}
        2x\equiv4\pmod 7,\\
        3x\equiv3\pmod 8,\\
        4x\equiv8\pmod 9
    \end{cases}
\]
\`e
\begin{enumerate}
    \item $1$
    \item $2$
    \item $3$
    \item $4$
    \item un numero diverso
\end{enumerate}
\end{samepage}

\begin{samepage}
\quesito{4}{3}
Il numero degli elementi invertibili nell'anello $\mathbb{Z}/1386\mathbb{Z}$,
espresso in base $5$, \`e
\begin{enumerate}
    \item $2420_5$
    \item $4560_5$
    \item $1220_5$
    \item $5500_5$
    \item $11300_5$
\end{enumerate}
\end{samepage}

\begin{samepage}
\quesito{5}{3}
Il resto della divisione per $73$ della potenza
$987654321^{123}$ \`e
\begin{enumerate}
    \item $2$
    \item $25$
    \item $9$
    \item $43$
    \item $1$
\end{enumerate}
\end{samepage}

\begin{samepage}
\quesito{6}{3}
Il numero di password di 10 simboli scelti con ripetizione dall'insieme
$\{a,e,i,o,u,w,x,y,z\}$ che contengono esattamente quattro vocali e le cui
consonanti compaiono in ordine alfabetico \`e
\begin{enumerate}
    \item $44100000$
    \item $11025000$
    \item un numero diverso
    \item $840000$
    \item $3652000$
\end{enumerate}
\end{samepage}

\begin{samepage}
\quesito{7}{3}
Sia $A=\{a,b,c,d\}$ e sia
\[
    R=\{(a,a),(b,b),(c,c),(d,d),(a,b),(b,c),(c,d)\}
\]
una relazione su $A$. La relazione $R$ \`e
\begin{enumerate}
    \item riflessiva, antisimmetrica e non transitiva
    \item riflessiva, non antisimmetrica e non transitiva
    \item riflessiva, non antisimmetrica e transitiva
    \item una relazione di ordine parziale
    \item simmetrica e transitiva
\end{enumerate}
\end{samepage}

\begin{samepage}
\quesito{8}{3}
Ponendo, per $a,b\in\mathbb{N}$,
\[
    a*b=\min\{a,b\},
\]
allora $(\mathbb{N},*)$ \`e
\begin{enumerate}
    \item un semigruppo, ma non un monoide
    \item un monoide, ma non un gruppo
    \item un gruppo
    \item non \`e un semigruppo
    \item un gruppo abeliano
\end{enumerate}
\end{samepage}

\begin{samepage}
\quesito{9}{2}
L'insieme
\[
    A=\{1,2,3,5,6,10,15\}
\]
ordinato rispetto alla relazione di divisibilit\`a
\`e
\begin{enumerate}
    \item un'algebra di Boole
    \item un reticolo distributivo
    \item non \`e un reticolo
    \item una catena
    \item un gruppo rispetto alla moltiplicazione
\end{enumerate}
\end{samepage}

\begin{samepage}
\quesito{10}{2}
Alla permutazione
\[
    \begin{pmatrix}
        1&2&3&4&5&6&7&8\\
        2&3&1&5&4&7&8&6
    \end{pmatrix}
\]
corrisponde la coppia (ordine, segno)
\begin{enumerate}
    \item $(6,\text{pari})$
    \item $(4,\text{dispari})$
    \item $(3,\text{pari})$
    \item $(6,\text{dispari})$
    \item $(8,\text{dispari})$
\end{enumerate}
\end{samepage}

\begin{samepage}
\quesito{11}{2}
Il numero delle stringhe ottenute permutando gli elementi dell'insieme
$A=\{a,b,c,d,e,f,g,h,i\}$ che non contengono la sottostringa \texttt{ab}
n\'e la sottostringa \texttt{cd} \`e
\begin{enumerate}
    \item $282240$
    \item $342120$
    \item un numero diverso
    \item $287280$
    \item $352800$
\end{enumerate}
\end{samepage}

\clearpage
\begin{center}
    {\Large\bfseries\color{examNavy} Scheda delle risposte}
\end{center}
Selezionare una casella per ogni quesito. Trascrivere il punteggio ottenuto
nel riquadro finale.

\vspace{0.8em}
\noindent
\begin{minipage}[t]{0.48\linewidth}
    \begin{tabular}{c*{5}{c}c}
        \toprule
        Quesito & A & B & C & D & E & Punti\\
        \midrule
        1 & $\square$ & $\square$ & $\square$ & $\square$ & $\square$ & \rule{0.45cm}{0.4pt}\\
        2 & $\square$ & $\square$ & $\square$ & $\square$ & $\square$ & \rule{0.45cm}{0.4pt}\\
        3 & $\square$ & $\square$ & $\square$ & $\square$ & $\square$ & \rule{0.45cm}{0.4pt}\\
        4 & $\square$ & $\square$ & $\square$ & $\square$ & $\square$ & \rule{0.45cm}{0.4pt}\\
        5 & $\square$ & $\square$ & $\square$ & $\square$ & $\square$ & \rule{0.45cm}{0.4pt}\\
        6 & $\square$ & $\square$ & $\square$ & $\square$ & $\square$ & \rule{0.45cm}{0.4pt}\\
        \bottomrule
    \end{tabular}
\end{minipage}\hfill
\begin{minipage}[t]{0.48\linewidth}
    \begin{tabular}{c*{5}{c}c}
        \toprule
        Quesito & A & B & C & D & E & Punti\\
        \midrule
        7 & $\square$ & $\square$ & $\square$ & $\square$ & $\square$ & \rule{0.45cm}{0.4pt}\\
        8 & $\square$ & $\square$ & $\square$ & $\square$ & $\square$ & \rule{0.45cm}{0.4pt}\\
        9 & $\square$ & $\square$ & $\square$ & $\square$ & $\square$ & \rule{0.45cm}{0.4pt}\\
        10 & $\square$ & $\square$ & $\square$ & $\square$ & $\square$ & \rule{0.45cm}{0.4pt}\\
        11 & $\square$ & $\square$ & $\square$ & $\square$ & $\square$ & \rule{0.45cm}{0.4pt}\\
        \bottomrule
    \end{tabular}
\end{minipage}

\vspace{2em}
\noindent
\colorbox{examPale}{%
    \parbox{\dimexpr\linewidth-2\fboxsep\relax}{%
        \textbf{Punteggio totale} \hfill
        \rule{1.5cm}{0.4pt}\,/\,30
        \hspace{1.5em}
        \textbf{Voto finale} \hfill
        \rule{1.5cm}{0.4pt}\,/\,30}}

\end{document}
